\section{Introduction}
An information state identifies histories that induce the same laws for
future experiments. A finite memory must do more than describe this
quotient: it must approximate the prediction directions actually
acquired, with the probability assigned to them by the experiment, and
must remain usable after every subsequent report. The dimension of the
quotient determines an asymptotic exponent in regular cases. It does
not determine the finite-resolution cost of directions that disappear
as the experiment or its prior degenerates.

This paper addresses these two questions separately. We first give an
experiment-level classification of the memory exponent for all positive
finite experiments whose likelihoods are polynomial in their commands.
The latent space need not be ordered, finite dimensional, or
semialgebraic, and the prior is arbitrary. We then determine the full
resolution profile for positive affine acquisition, including covariance
rank loss. A necessary and sufficient criterion characterizes
prior-uniform attainment in the square pairing regime; a product-cone
alternative covers arbitrary rectangular pairings. We also realize
arbitrary small covariance matrices in a fixed experiment and classify
their analytic degenerations jointly with the memory budget. Two multi-step
classifications identify the finer geometry of additive exponent
collisions and of vanishing circular contrast. In each resolution law,
the scales concern the interaction of acquisition and future observation,
not the future test space alone.

\subsection{An experiment-level rank and its prior strata}
Let $X$ be a compact metric space, let $\mu$ be a probability on $X$,
and fix a finite horizon $N$. At time $i$ the command belongs to a
finite-dimensional cube, the report alphabet is finite, and
\[
 L_{i,y}(u,x)=\sum_{\alpha\in F_{i,y}}u^\alpha c_{i,y,\alpha}(x),
 \qquad c_{i,y,\alpha}\in C(X),\qquad L_{i,y}\ge\kappa>0,
 \qquad\sum_yL_{i,y}=1.
\]
The polynomial condition concerns $u$, not $x$. In particular, finite
detectors with affine retention gates belong to this class for arbitrary
continuous detector cells. Future word likelihoods span a
finite-dimensional test space. Choose a finite spanning list of actual
queries, each with a fixed positive selection probability, and write
$p_n(h)$ for the weighted vector of their conditional success
probabilities after a history of length $n$.

For each past report word $a$, let $P_a$ be its likelihood product as a
function of the commands, set $Z_a=\mu P_a$, and put
\[
 Y_a=(Z_a,Z_ap_n)^t,\qquad
 d_n=\max_a\rank_{\R(\boldsymbol u)}[Y_a,DY_a]-1.
\]
All entries are command polynomials with finitely many actual prior
moments as coefficients. The evidence column is adjoined before
normalization; it is not counted as an acquired prediction direction.
Theorem~\ref{thm:v18-rank-classification} proves, whenever $d_n>0$,
\[
 R^{\rm av}_{M,n}\asymp R^{\rm max}_{M,n}\asymp M^{-2/d_n},
 \qquad
 \mathfrak R^{\rm av}_{M,N}\asymp
 \mathfrak R^{\rm max}_{M,N}\asymp M^{-2/d_*},
 \quad d_* =\max_{1\le n<N}d_n>0.
\]
Here $R$ is checkpoint regret and $\mathfrak R$ is the largest
checkpoint regret of one causal filter. The average criterion uses
independent uniform commands and includes report evidence; the maximum
criterion ranges over all admitted histories. A zero rank gives a
finite reachable set, not necessarily a singleton. If all checkpoint
ranks vanish, a finite exact causal filter exists.

This formula classifies each specified experiment without assuming an
accessible flag or an ordered-scale representation of its risk.
Theorem~\ref{thm:v18-moment-strata} further identifies every locus
$d_n\le k$ by the vanishing of finitely many polynomials in the prior
moments: they are the command coefficients of the $(k+2)$-minors of the
augmented matrices. Full-support priors fill the interior of the finite
moment body. Simultaneous maximal ranks hold off a proper algebraic
set there, and can be reached by arbitrarily small bounded-density
perturbations of any full-support prior. Exceptional priors are not
excluded from the classification; their actual ranks enter the formula.
Uniform constants hold on compact constant-rank moment sets. A rank
formula by itself does not assert uniformity across changing ranks.

\subsection{A complete profile from acquisition--query covariance}
The first full-class resolution theorem applies to positive affine
acquisition. Let $f\in C(X)^b$ with
$\sup_x\sum_i|f_i(x)|\le a<1$, choose a uniform command
$u\in[-1,1]^b$, and observe a report $y\in\{-1,1\}$ with likelihood
\[
 L_y(u,x)=\tfrac12(1+y u^tf(x)).
\]
For any finite menu of future query probabilities, let $h$ be its vector
weighted by the square roots of the query probabilities and define
\[
 C_\mu=\Cov_\mu(h,f).
\]
The physical command cube fixes the input units. No singular whitening
is performed. If $s_1\ge\cdots\ge s_d\ge0$ are the singular values
of $C_\mu$, Theorem~\ref{thm:v18-affine-classification} proves
\[
 R_M^{\rm av}\asymp R_M^{\rm max}\asymp
 \max_{1\le\ell\le d}
      \left(\frac{s_1\cdots s_\ell}{M}\right)^{2/\ell}.
\]
The constants depend only on $a,b$ and the number of queries. They are
uniform over priors, query weights, multiplicities, and all rank losses.
The assertion includes the one-label exact case $C_\mu=0$.

The proof derives this geometry before addressing quantization. With
$v=yu$ and $m=\mu f$, the exact prediction vector is
\[
 p=\mu h+C_\mu\theta,\qquad
 \theta=\frac{v}{1+m^tv}.
\]
The law of $v$, including both reports and their evidence, has density
$(1+m^tv)/2^b$. Its projective image contains a fixed ball, is contained
in another fixed ball, and has a uniformly positive density on the
inner ball. Thus the same covariance image controls both global
approximation and unconditional acquisition mass. This is not an
inference from an ambient observation matrix. As a concrete interior
rank change, acquisition through $x$ and prediction through $x^2$ under
$d\mu_t=(1+tx)\,dx/2$ on $[-1,1]$ have covariance proportional to
$t$. Proposition~\ref{prop:v18-interior-example} obtains the full
uniform $t^2M^{-2}$ law, including zero at $t=0$, while every prior
still has full support.

Theorem~\ref{thm:v18-universal-pairing} addresses a different
quantifier. At an interior positive history let
$\widetilde T=\operatorname{span}\{P\}+\operatorname{im}DP$ and
$\Phi=\operatorname{span}\{1,\phi_1,\ldots,\phi_p\}$. When both
spaces have dimension $p+1$, normalized acquisition has rank $p$ for
\emph{every} full-support prior if and only if
\[
 \det[v_i(x_k)]_{i,k=0}^p\det[\phi_j(x_k)]_{j,k=0}^p
\]
has one weak sign on $X^{p+1}$ and is nonzero somewhere. Here $v$ is
any basis of $\widetilde T$. Necessity follows by approximating atomic
probabilities by full-support probabilities; sufficiency follows from
determinant integration. Separate Chebyshev properties are sufficient
instances, not necessary hypotheses. In particular,
$\widetilde T=P\Phi$ gives prior-uniform full rank on arbitrary compact
latent spaces. The square hypothesis is explicit; no reduction of every
rectangular pairing to this criterion is asserted.

\subsection{Rectangular attainment and realizable degenerations}
For finite-dimensional $E,F\subset C(X;\R)$ with $\dim F\le\dim E$,
Theorem~\ref{thm:v19-rectangular} characterizes full pairing rank
under every full-support prior by
\[
 0\ne f\in F\quad\Longrightarrow\quad
 \text{there exists }g\in E\text{ with }fg\ge0,\quad fg\not\equiv0.
\]
Theorem~\ref{thm:v19-rank-alternative} additionally identifies all
possible kernel subspaces through relative interiors of product-moment
bodies and constructs full-support priors realizing the extremal rank.
Adjoining the positive history product before normalization converts
pairing rank into actual prediction rank. This yields the unconditional
mass conclusion without a square reduction. The five-point experiment
in Proposition~\ref{prop:v19-pentagon} satisfies the rectangular
criterion but has no fixed universally nonsingular two-dimensional
tangent subspace containing evidence.

A second conclusion concerns rank degeneration inside a dominated
prior family. For each pair $b,q$, one fixed positive acquisition and
one fixed physical query menu admit priors $\mu_A$ with
\[
 \tfrac12\mu_0\le\mu_A\le\tfrac32\mu_0,\qquad
 C_{\mu_A}=\frac{A}{8b\sqrt q},\qquad
 \sum_{j,i}|A_{ji}|\le\tfrac12.
\]
Thus arbitrary small matrix germs, and not only a scalar cancellation,
are realized by genuine priors in one experiment
(Theorem~\ref{thm:v19-covariance-realization}). Along an analytic
arc let $A_\ell$ be the least vanishing order among its $\ell$-minors.
Theorem~\ref{thm:v19-analytic-memory} gives, uniformly in small
$t>0$ and all integer $M\ge1$,
\[
 R_{M,1}(t)\asymp
     \max_{1\le\ell\le r}t^{2A_\ell/\ell}M^{-2/\ell}.
\]
The differences $\alpha_\ell=A_\ell-A_{\ell-1}$ are ordered
nonnegative integers. Every such finite list occurs under the same
physical bounds, and the full checkpoint budget curve determines the
$A_\ell$ operationally. In particular the phase exponent along
$M=\lceil t^{-\gamma}\rceil$ is
$\min_\ell 2(A_\ell+\gamma)/\ell$. This is a complete arcwise
classification for affine acquisition, not a simultaneous normal form
for all polynomial history maps in several parameters.

\subsection{Uniform multi-step geometry through exponent collisions}
For the monomial experiment write
\[
 A=\{0=a_0<a_1<\cdots<a_{r-1}=D\},\qquad
 W_A=\operatorname{span}\{t^a:a\in A\},\qquad h_A(m)=|mA|,
\]
where $mA$ consists of sums of exactly $m$ members of $A$, including
zero entries. A positive finite detector spans $W_A$. A command in
$[\eta,1-\eta]^J$, with $0<\eta<1/2$, retains or rejects each cell;
rejection is a report, and every trial is counted. All report
likelihoods are at least $\kappa>0$. The exact information statements
hold on $I=[l,u]$, $0\le l<u<\infty$, with a full-support prior.
For the uniform collision statements take $I=[0,1]$, fix $N\ge2$,
and let the one-step exponents vary over a compact subset $K$ of the
strictly ordered chamber. One fixed full-column-rank coefficient matrix
specifies the positive detector on $K$, and $\mu$ is fixed with full
support; it need not have a density.

Choose $H\ge\max\{1,N\max_K a_{r-1}\}$. Keep every nonconstant
formal future label $\alpha\in\N_0^r$, $|\alpha|=m$, even when
exponents coincide. Put
\[
 q_m=\binom{m+r-1}{r-1}-1,\qquad x_\alpha=H^{-1}\alpha\cdot a,
 \qquad
 \mathcal V_{m,\ell}(a)=\max_{|F|=\ell}
       \prod_{\{\alpha,\beta\}\subset F}|x_\alpha-x_\beta|.
\]
Thus $\mathcal V_{m,1}=1$; a zero product remains zero. Set
\[
 \begin{split}
 p_{n,m}&=\min\{n(r-1),q_m\},\\
 \Psi_{n,m}(M,a)&=\max_{1\le\ell\le p_{n,m}}
       (\mathcal V_{m,\ell}(a)/M)^{2/\ell},\\
 \Xi_N(M,a)&=\max_{1\le n<N}\Psi_{n,N-n}(M,a).
 \end{split}
\]
Endpoint profiles are zero. The checkpoint query is a uniformly selected
member of a fixed finite spanning menu of physical $m$-trial products,
selected independently after the memory index is formed. Only the
selected query is executed. Commands during acquisition are independent
uniform commands for the unconditional criterion. The resource and
worst-history and causal criteria are those defined above and in
Definition~\ref{def:finite-state}.

\begin{theorem}[Intrinsic attainable resolution and causal memory]\label{thm:resolution-main}
Under these hypotheses there are constants $0<c\le C<\infty$,
depending only on the fixed experiment, $K$, $N$ and $\mu$, with the
following properties for every known calibration $a\in K$ and every
integer $M\ge1$.

At each checkpoint $n+m=N$, the optimal $M$-message squared-prediction
regret is between $c\Psi_{n,m}(M,a)$ and $C\Psi_{n,m}(M,a)$,
both unconditionally under independent uniform acquisition commands and
in the minimax sense over admitted histories. The query is sampled
independently after the index is formed, from the ordered products of
one fixed attainable failure-factor basis; all its $m$ trials are counted.

There exists one deterministic stage-compatible transducer, allowed to
depend on $a$ and $M$, storing at most $M$ persistent labels after
every report, whose regret at every checkpoint and on every history is
at most $C\Xi_N(M,a)$. Conversely every such transducer has maximum
unconditional checkpoint regret at least $c\Xi_N(M,a)$ under the
same common exploration law. Independent coding randomization does not
remove the lower bound. The clock and fixed real-valued program are
read-only as in Definition~\ref{def:finite-state}.

The constants are uniform across all additive collision strata of $K$.
No calibration-blind codebook or uniformity over all full-support priors
is part of the assertion.
\end{theorem}


The exact continuous-state dimension is separately
$d_A(n,m)=\min\{n(r-1),|mA|-1\}$, as stated in
Theorem~\ref{thm:main}. The finite-resolution formula specifies how
these directions disappear. Its proof joins three assertions not
contained in an ambient spectral estimate. Actual interior failure
factors have a separated binomial product tangent. Pairing it with
complete Newton--Hermite prefixes loses exactly the evidence direction
under normalization and supplies a positive-history minorization.
The entire posterior image has bounded semialgebraic format, so its
thin-rectangle covering bound truncates at the same attainable
dimension. Finally, raw-moment updates have a Lipschitz recurrence
without inverse collision gaps, and quantization of reachable
representatives realizes one causal filter. These are established in
Lemmas~\ref{lem:newton-attainment} and \ref{lem:tame-rectangle} and
Theorems~\ref{thm:intrinsic-checkpoint} and
\ref{thm:intrinsic-streaming}.

\subsection{Contrast, inversion, and unresolved prior information}
A circular experiment with Haar prior has the positive cells
$\tfrac14(1\pm\tau\cos\theta)$ and
$\tfrac14(1\pm\tau\sin\theta)$. Its four lookup probabilities vary
continuously in $[\eta,1-\eta]^4$. For every fixed horizon and all
sufficiently small known $\tau\ge0$, the checkpoint law is
\[
 H_{n,m}(M,\tau)=
   \max_{1\le j\le\min(n,m)}\tau^{2(j+1)}M^{-1/j}.
\]
The $j$th harmonic is attenuated by $\tau^j$ when acquired and again
by $\tau^j$ when queried. The real physical scales are therefore
$\tau^2,\tau^2,\tau^4,\tau^4,\ldots$, not just a count of Fourier
coefficients. Theorem~\ref{thm:circle-resolution} proves a uniform
causal realization through weighted Fourier updates without inverse
contrast. The contrast interval and constants belong to the fixed
horizon, not to an unbounded-horizon assertion.

Within these established profile families the complete checkpoint curve
has an inverse. For $R(M)$ define
\[
 \mathcal I_\ell(R)=\inf_{M\in\N,\,M\ge1}M R(M)^{\ell/2}.
\]
Theorem~\ref{thm:operational-reconstruction} and
Corollary~\ref{cor:v18-affine-invariant} identify this transform with
the initial scale products up to uniform constants. Integer budgets
require a factor-two sandwich; the underlying real-budget identity is
exact. Repeated scales and zero products are included. This inverse
uses the entire unbounded budget curve, not finitely observed errors,
and does not recover a latent model or exact optimal constants. Nor is
it applied to a maximum over checkpoints or to a curve with a positive
uncertainty floor.

Prior uncertainty has its own geometry. For monomials, the local
ambiguity body after a positive history has successive width orders
$\varepsilon d_1,\ldots,\varepsilon d_{q_m}$, where $d_i$ are the
future Leja scales. There is no past-dimension truncation of this
ambiguity body. After an exact prefix of $k$ raw future moments,
the remaining squared radius is $\asymp\varepsilon^2d_{k+1}^2$.
The normalization of prior tilts is explicit in
Proposition~\ref{prop:ambiguity-ellipsoid}. For an imperfect common
moment name of error $\delta$, the separate dominated-family and
interior-slack hypotheses of Theorem~\ref{thm:sharp-common-moments}
give the sharp joint law $\Xi_N(M,a)+\delta^2$. The lower bound
uses fixed genuine prior alternatives and overlapping history laws,
not alternatives chosen after seeing a history.

\subsection{Method, attribution, and organization}
The analytic inputs are semialgebraic dimension and metric entropy,
determinant integration, total positivity, and divided differences
\cite{BochnakCosteRoy,ForresterAndreief,Pinkus,deBoor,
YomdinComte,ZhangKileel,ComteHalupczok}. They are used as classical
inputs, not separate novelty claims. The finite square-pairing
necessity criterion is the strong-compatibility result of
Banaji--Pantea \cite{BanajiPantea2016}; the finite rectangular
sign-vector alternative is present in M\"uller et al.\
\cite{MullerEtAl2016}. Analytic invariant orders and their
singular-value interpretation are classical as well
\cite{KavehMakhnatch}. Section~\ref{sec:v19-comparison}
records the exact specializations, rather than attributing only the
underlying determinant identity. The structural result identifies
their consequences for an explicitly defined experiment and its actual
prior. Its causal part requires that future test spaces close under
multiplication by the next likelihood; its lower part requires the
report-word evidence and integration in the unused command coordinates.
The affine classification proves a uniform global image and mass
comparison directly from Bayes' formula. The multi-step classifications
require the additional model-specific acquisition arguments described
above. A dimension theorem is not substituted for any of these uniform
profile statements.

The manuscript develops the structural classification, affine profile,
their analytic degeneration, and the square and rectangular pairing
criteria first, then the complete monomial and
circular classifications, their operational inverse, and prior
uncertainty. The observation algebra, confluent specializations,
sequential decision results, and all numerical constructions appear
with their full proofs in the appendices. In the exact-real model,
calibration, clock, and fixed program are read-only; no uncharged
history or command-generation seed is allowed. The separate finite
implementation theorem, Theorem~\ref{thm:effective-main}, supplies
integer transition tables and dyadic outputs under its explicit
numerical interface. Its advice, program length, transient workspace,
and persistent labels are distinct resources. Sharpness is asserted
for persistent memory; sufficient bounds are supplied for the other
resources. The intrinsic results do not assume computability of
unsupplied real prior moments. Detailed comparisons with finite
control approximations and information-state methods are retained in
Appendix~\ref{sec:literature-comparison}.
